\section{Introduction}

Smoothed-particle hydrodynamics (SPH) is a particle-based Lagrangian method used to simulate multidimensional fluids, most commonly employed in astrophysics, cosmology and computational fluid-dynamics. First developed by Gingold and Monaghan\cite{GingoldMonaghan1977} and Lucy\cite{Lucy1977} in 1977, it can outshine traditional grid-based Eulerian codes in some key areas such as its natural conservation properties \cite{Rosswog2009}, the handling of free surfaces and fluid interfaces \cite{He2024}, adaptive resolution \cite{Price2010} and it integrates well with fast gravity solvers for self-gravitating systems, which is particularly useful in astrophysical simulations. 

SPH-EXA is a high-performance \texttt{C++} code developed by the University of Basel for performing such SPH simulations on Exascale computing systems \cite{Cavelan2020}. It is parallelized with \texttt{MPI}, \texttt{OpenMP}, \texttt{CUDA}, and \texttt{HIP} to enable calculations with billions of particles running on thousands of hybrid CPU-GPU nodes in parallel. 

In most, if not all astrophysical problems, the influence of magnetic fields on plasma behavior is crucial in accurately describing them.\todo{concrete example} Therefore, many modern codes include a magnetohydrodynamic (MHD) extension for SPH (SPMHD). The primary goal of this project is to present a working and properly scaling implementation of SPMHD in SPH-EXA. The scalability of SPH-EXA will hopefully allow future simulations of large-scale MHD problems, which were previously limited by resolution constraints. 

In the following chapters we first briefly touch on the fundamentals of the SPH discretizations and how they are implemented in SPH-EXA using an integral approach to derivation. This will provide the necessary basis and context to then introduce the newly added magnetic field calculations. We base our implementation on other state-of-the-art SPMHD codes, namely \textsc{Phantom} \cite{Price2018} and the MHD extension to \textsc{Gasoline2} \cite{Wadsley2017, WissingShen2020}. 

We will give special focus to the treatment of discontinuities through artificial resistivity. Different approaches that emerged throughout the years are discussed. Then, we present a novel approach to limit unwanted dissipation, using slope-limited linear reconstruction of the local magnetic field, based on the same method recently proposed for artificial viscosity in \cite{GarciaSenzCabezon2026}. 

We then conduct a series of numerical tests to verify that our implementation produces results comparable to other SPMHD codes. There we also study the impact of different AR schemes, comparing our new design to a baseline implementation based on \cite{WissingShen2020}. As a final test we simulate subsonic magnetic turbulence at $N=1600^3$ particles as a proof of concept that SPH-EXA is now properly equipped to take on real astrophysical MHD problems beyond standardized tests. 

In chapter \ref{sec:performance} we study performance in a strong and weak scaling analysis, ensuring that the MHD extension scales together with the hydrodynamic parts of SPH-EXA. We also quantify the exact overhead on runtime in a per-kernel analysis. 

Finally, the key findings are discussed and next steps are outlined. 

Throughout this thesis paper we will heavily rely on Wissing \& Shen (2020) \cite{WissingShen2020} and Price et al. (2018) \cite{Price2018} as our primary reference both regarding the numerical method and as direct comparison for the tests in our verification suite. 

It is important at this point to properly credit Helena Schmidt (University of Basel), who wrote considerable parts of the current MHD module in SPH-EXA. The bulk work of this thesis project could therefore be focused on updating, validating, extending, optimizing and formalizing this implementation, rather than building it from scratch. 